We introduced HeadIG, a DLLM-aligned head-attribution framework that defines head importance through the masked reconstruction objective rather than autoregressive next-token prediction. By inserting differentiable gates on attention-head outputs and integrating gradients with respect to these gates, HeadIG turns head attribution in diffusion language models into a principled continuous attribution problem.

We further proposed PolyHeadIG, a multi-path variant built on stochastic threshold paths. This construction alleviates the synchronization bias of single-path diagonal IG and provides a more expressive attribution mechanism over asynchronous gate configurations. In addition to the method itself, we established an evaluation protocol that compares HeadIG against AttnLRP and Shapley-value baselines under two causal validations: attribution-guided adaptive sparse inference and importance-based head masking and pruning.

More broadly, the framework developed in this paper aims to make head attribution in DLLMs both methodologically aligned and practically actionable. Beyond interpretability, these attribution scores can support downstream efficiency mechanisms and provide a tool for studying how attention heads contribute across denoising regimes, tasks, and model families.

Several directions remain for future work. An important next step is to study richer path families and lower-variance estimators for attribution in high-dimensional gate spaces. It is also natural to extend the framework beyond attention heads to other internal components in diffusion language models, and to investigate whether attribution signals transfer across models, tasks, and inference regimes.
